% Generated draft: requires mathematical and source review.
\documentclass[11pt]{article}
\usepackage[margin=25mm]{geometry}
\usepackage{fontspec}
\setmainfont{Latin Modern Roman}
\usepackage{amsmath,amssymb,mathtools,unicode-math}
\setmathfont{Latin Modern Math}
\usepackage{enumitem,needspace,xurl,hyperref}
\setlength{\parindent}{0pt}
\setlength{\parskip}{5pt}
\setlength{\emergencystretch}{3em}
\newcommand{\N}{\mathbb{N}}
\newcommand{\Z}{\mathbb{Z}}
\newcommand{\Q}{\mathbb{Q}}
\newcommand{\R}{\mathbb{R}}
\newcommand{\C}{\mathbb{C}}
\newcommand{\F}{\mathbb{F}}
\newcommand{\A}{\mathbb{A}}
\newcommand{\PP}{\mathbb{P}}
\newcommand{\E}{\mathbb{E}}
\newcommand{\eps}{\varepsilon}
\newcommand{\dd}{\,\mathrm{d}}
\begin{document}
\section*{UnsolvedMath ID: 800901}
\textbf{Source problem code:} catalogue entry
\subsection{Definitions and mathematical statement}
A quasi-plurisubharmonic function is a function locally expressible as the sum of a plurisubharmonic function and a smooth real-valued function, with the usual exclusion of functions identically \(-\infty\) on a connected component. Write \(dd^c\) using a fixed consistent normalization, and let \(\nu(R,x)\) denote the Lelong number of a positive closed \((1,1)\)-current \(R\) at \(x\).

The supplied excerpt uses an ambient complex manifold \(X\), a form \(\omega\), and a quasi-plurisubharmonic function \(\varphi\), but does not state their full hypotheses. To express its proposed conclusion, suppose the datum admits the positive closed \((1,1)\)-current
\[
T=\omega+dd^c\varphi.
\]
For such a datum, the requested desingularization property is the following:
\[
\begin{gathered}
\text{For every }\varepsilon>0\text{ there exist a finite composition of blow-ups}\quad
\pi_\varepsilon:X_\varepsilon\longrightarrow X,\\
D_\varepsilon=\sum_{j=1}^{m_\varepsilon}a_{\varepsilon,j}D_{\varepsilon,j},
\qquad a_{\varepsilon,j}>0,
\qquad\text{and a positive closed }(1,1)\text{-current }R_\varepsilon,\\
\pi_\varepsilon^*T=[D_\varepsilon]+R_\varepsilon,
\qquad
\nu(R_\varepsilon,x)<\varepsilon\quad\text{for every }x\in X_\varepsilon.
\end{gathered}
\]
Here \(X_\varepsilon\) is smooth, the \(D_{\varepsilon,j}\) are distinct irreducible hypersurfaces, and their union has simple normal crossings: locally it is a union of coordinate hyperplanes. The sum may be empty. The current \([D_\varepsilon]\) means \(\sum_j a_{\varepsilon,j}[D_{\varepsilon,j}]\), where \([D_{\varepsilon,j}]\) is the current of integration over that hypersurface.

\textbf{Question 32, attributed in the source to Favre--Guedj--Jonsson.} Does this property hold in every complex dimension \(n=\dim_{\C}X\geq3\), for all data in the intended ambient setting of the source?

The supplied excerpt reports that the property holds when \(n\leq2\), citing Blel--Mimouni [BM05], Guedj [G05b], and Favre--Jonsson [FJ05]. This is a supplied background assertion. The excerpt does not specify the ambient class of \(X\), the assumptions on \(\omega\), or the permitted blow-up centers. Consequently, the quantified desingularization property above is explicit, but the class of initial data over which the question is intended to range remains unresolved; in particular, no assertion for all complex manifolds is being substituted for the source question.
\subsection{Short English statement}
The question asks whether, in complex dimension at least three, finitely many blow-ups can separate the singularities of the given current into a divisor with simple normal crossings and a positive remainder whose Lelong number is strictly less than any prescribed \(\varepsilon>0\) at every point. The supplied excerpt reports this in dimensions at most two, but omits ambient hypotheses needed to fix the exact higher-dimensional statement.
\subsection{Sources}
\begin{itemize}
\item Local registry: AMR Open Problem Lists, permanent UnsolvedMath ID \(800901\), exact source problem code \texttt{catalogue entry}, title \emph{Open Problems in Pluripotential Theory --- Question 32}. The supplied registry labels the problem \texttt{open}; this is a supplied status label, not an independently verified current assessment.
\item Dinew, Guedj, Zeriahi, \emph{Open Problems in Pluripotential Theory} (2015), Question 32. Supplied source URL: \url{https://arxiv.org/abs/1511.00705}. The record reports extraction by \texttt{arxiv\_latex} and an access date of 2026-10-16; these are supplied provenance details, not a claim of independent access or verification.
\item The question itself carries the attribution Favre--Guedj--Jonsson. The registry separately lists Dinew, Guedj, Zeriahi as proposers, with proposed year 2015.
\item Supplied references for the dimension-at-most-two background assertion: Blel--Mimouni [BM05], the second author of the source list, Guedj [G05b], and Favre--Jonsson [FJ05]. The supplied record does not include full bibliographic entries for these labels.
\end{itemize}
\end{document}
